\documentclass[11pt,a4paper]{article}

\usepackage[utf8]{inputenc}
\usepackage[T1]{fontenc}
\usepackage{geometry}
\geometry{
    a4paper,
    left=20mm,
    right=20mm,
    top=20mm,
    bottom=20mm
}
\usepackage{titlesec}
\usepackage{enumitem}
\usepackage{xcolor}
\usepackage{hyperref}
\usepackage{parskip}

% Formatting for sections
\titleformat{\section}
  {\Large\bfseries\color{blue!60!black}}
  {\thesection.}
  {1em}
  {}

\titleformat{\subsection}
  {\large\bfseries\color{blue!80!black}}
  {}
  {0em}
  {}

\hypersetup{
    colorlinks=true,
    linkcolor=blue,
    urlcolor=blue
}

\begin{document}

\begin{center}
    {\Huge \textbf{VisionGuard AI}} \\[0.5cm]
    {\LARGE \textbf{Presentation \& Viva Cheat Sheet}} \\[0.5cm]
    {\large Muhammad Ijlal Asif (F24607072) \quad|\quad Syed Maaz Shah (F24607091)}
\end{center}

\vspace{0.5cm}
\hrule
\vspace{0.5cm}

\textit{Use this quick guide to answer any questions the teacher might ask during your presentation or viva.}

\section{Core Concept (What did we build?)}
\begin{itemize}
    \item \textbf{What is it?} An AI-based web application to detect Glaucoma from retinal fundus images.
    \item \textbf{Dataset Used:} ACRIMA Dataset (705 images: 396 Glaucoma, 309 Normal).
    \item \textbf{Why is it important?} Glaucoma has no early symptoms. Our tool helps doctors screen patients quickly and cheaply, especially in rural areas where eye specialists are rare.
\end{itemize}

\section{The Models (Teacher vs. Student)}
\begin{itemize}
    \item \textbf{Why two models?} We used a technique called \textbf{Knowledge Distillation}.
    \item \textbf{The Teacher (ResNet-50):} A massive, slow, but very accurate model. It learns the complex patterns perfectly but is too heavy (153 MB) to run smoothly on a web app.
    \item \textbf{The Student (EfficientNet-B0):} A very small, fast model (only 24 MB).
    \item \textbf{How it works:} Instead of training the Student from scratch, we force the Student to copy the ``knowledge'' of the Teacher. The Student learns to mimic the Teacher's decisions, giving us the accuracy of a big model with the speed of a small one.
\end{itemize}

\section{Multi-Task Learning (Doing two things at once)}
\begin{itemize}
    \item \textbf{Old way (Assignments 1-3):} The model only did one thing: output ``Normal'' or ``Glaucoma'' (Classification).
    \item \textbf{New way (Final Project):} Our network does \textit{two} things simultaneously:
    \begin{enumerate}
        \item \textbf{Classification:} Diagnoses Normal or Glaucoma.
        \item \textbf{Segmentation:} Physically maps out the Optic Disc and Optic Cup pixel-by-pixel to calculate the CDR (Cup-to-Disc Ratio).
    \end{enumerate}
    \item \textbf{Why?} Making the AI learn the physical structure of the eye (segmentation) forces it to be smarter and prevents it from making random guesses.
\end{itemize}

\section{CBAM Attention (The ``Focus'' Mechanism)}
\begin{itemize}
    \item \textbf{What is it?} Convolutional Block Attention Module.
    \item \textbf{What does it do?} It acts like a spotlight for the AI. It tells the network \textit{what} to look at (Channel Attention) and \textit{where} to look (Spatial Attention).
    \item \textbf{Why use it?} It forces the AI to ignore background noise (like blood vessels or camera flash) and focus 100\% on the Optic Nerve Head, exactly like a real doctor does.
\end{itemize}

\section{Explainable AI (XAI) - Proving it's not ``Fake''}
\begin{itemize}
    \item \textbf{The Problem:} Doctors don't trust ``Black Box'' AI that just spits out an answer without explaining why.
    \item \textbf{Our Solution:} We added two XAI methods so the doctor can literally see what the AI is thinking.
    \begin{enumerate}
        \item \textbf{Grad-CAM:} Creates a colored heat map (coarse localization) showing the general area the AI focused on.
        \item \textbf{Integrated Gradients (IG):} A more advanced, pixel-level map showing exactly which individual pixels influenced the decision.
    \end{enumerate}
\end{itemize}

\section{Real-World ``Authenticity'' Features (Crucial!)}
\textit{If the teacher asks why the probability is never exactly 100\% or 0\% anymore, explain these features:}
\begin{itemize}
    \item \textbf{Temperature Scaling:} Deep learning models are usually overconfident (spitting out 99.9\% or 0.1\%). We added a math formula (Temperature = 2.5) to ``soften'' the predictions. This makes the percentages realistic (e.g., 85\% instead of 99.9\%), proving it's an authentic AI, not a hardcoded fake.
    \item \textbf{Out-of-Distribution (OOD) Detection:} We calculate the ``Entropy'' (confusion level) of the image. If someone uploads a picture of a cat or a car, the Entropy spikes, and the system rejects the image with a warning instead of trying to diagnose it.
    \item \textbf{Test-Time Augmentation (TTA):} When a user uploads an image, the backend flips it horizontally, flips it vertically, and runs the original. It averages the score of all three. This makes the final prediction much more stable and reliable.
    \item \textbf{Adjustable Threshold:} The slider on the UI lets the doctor decide how strict the AI should be. Lowering the threshold catches more Glaucoma cases (good for screening), while raising it prevents false alarms.
\end{itemize}

\section{Tech Stack (What technologies were used?)}
\begin{itemize}
    \item \textbf{Deep Learning:} PyTorch, Torchvision.
    \item \textbf{Backend:} FastAPI (Python) -- Handles the heavy lifting, AI processing, and image generation.
    \item \textbf{Frontend:} React (Vite) -- Creates the beautiful, fast, glassmorphism UI.
    \item \textbf{Why FastAPI + React instead of Streamlit?} Streamlit was good for a prototype (Assignment 3), but FastAPI + React is how actual, professional software is built in the real industry. It separates the front and back for better speed and scalability.
\end{itemize}

\end{document}
